% JAIR requires a short title for the page headers, since the full one is
% three lines. The line breaks the full title used under acmart are gone:
% the AuthorKit says not to insert breaks in a title, and jair.cls wraps it.
\title[A Rubric-Based Survey of Inter-Agent Protocols]{A Rubric-Based Survey
  of Inter-Agent Protocols: Industry Standards, Security Posture, Open
  Problems, and Three Co-Designed Protocol Extensions}

\author{Suranjan Goswami}
\authornote{Corresponding Author.}
\affiliation{%
  \institution{Nasiko Labs}
  \country{India}
}
\email{suranjan_arcade@yahoo.co.in}

\author{Manu Ghulyani}
\affiliation{%
  \institution{Nasiko Labs}
  \country{India}
}
\email{manughulyani@gmail.com}

\author{Karan Bharadwaj}
\affiliation{%
  \institution{Nasiko Labs}
  \country{USA}
}
\email{karan@nasiko.com}

% Details taken from the public record only, pending his own confirmation.
% Name string, affiliation and email are exactly as he self-published them on
% the title page of arXiv:2601.14567v2, the agent:// scheme paper cited here
% as \cite{agent_uri_scheme}. That page reads "Roland R. Rodriguez, Jr. /
% Independent Researcher / rrrodzilla@proton.me", and its footnote 1 names
% the same address as the contact point. No country appears on that page;
% United States comes from his publicly listed Govcraft office in Corpus
% Christi, Texas (govcraft.ai). The suffix comma is braced so acmart does
% not read it as an author separator.
\author{Roland R. Rodriguez, {Jr.}}
\affiliation{%
  \institution{Independent Researcher}
  \country{United States}
}
\email{rrrodzilla@proton.me}

\author{Faouzi El Yagoubi}
\affiliation{%
  \institution{Polytechnique Montr\'eal}
  \city{Montr\'eal}
  \state{Qu\'ebec}
  \country{Canada}
}
\email{faouzi.el-yagoubi@polymtl.ca}

\author{Alejandro Quintero}
\affiliation{%
  \institution{Polytechnique Montr\'eal}
  \city{Montr\'eal}
  \state{Qu\'ebec}
  \country{Canada}
}
\email{alejandro.quintero@polymtl.ca}

\author{Ranwa Al Mallah}
\affiliation{%
  \institution{Polytechnique Montr\'eal}
  \city{Montr\'eal}
  \state{Qu\'ebec}
  \country{Canada}
}
\email{ranwa.al-mallah@polymtl.ca}

% JAIR asks for a short author list for the running head, and strongly
% recommends an \orcid per author. Neither the Nasiko Labs authors nor
% Rodriguez has supplied an ORCID, and the Polytechnique Montreal authors
% joined too recently to have been asked, so no \orcid lines appear above.
% They are not invented here: a wrong ORCID resolves to a real stranger.
% Collect all seven before camera-ready; the README tracks this.
\renewcommand{\shortauthors}{Goswami et al.}

% =====================================================================
% Reserved author slot (1 remaining of the 3 originally held).
% Slots six and seven went to Alejandro Quintero and Ranwa Al Mallah
% (Polytechnique Montreal), named as co-authors in Faouzi El Yagoubi's
% reply of 2026-09-09.
%
% Kept commented so the circulating draft's title page does not show
% placeholder names. The block needs all four fields, or acmart emits
% a missing-metadata warning.
% =====================================================================
% \author{Author Eight}
% \affiliation{%
%   \institution{TODO: affiliation}
%   \country{TODO: country}
% }
% \email{TODO@example.org}

% JAIR uses a structured abstract: Background, Objectives, Methods, Results,
% Conclusions, each introduced by a bold run-in label. The content is the same
% claim set as the prose abstract this replaced, redistributed under the five
% headings, with the honesty disclosures (single rater, no inter-rater
% statistic, sketches not standards) kept where a reader meets them early.
% Note that the arXiv abstract in manuscript/arxiv-abstract.txt is a third,
% independent text under a 1920-character limit. Changing a claim here means
% changing it there too; see manuscript/arxiv-abstract-NOTES.md.
\begin{abstract}
{\bf Background:}
The 2024 to 2026 period has produced an explosion of inter-agent protocols.
Google A2A, Anthropic MCP, IBM ACP, Cisco AGNTCY/SLIM, ANP, Coral, LOKA,
ACNBP, AP2, Visa Trusted Agent Protocol, ERC-8004, and UCP each address
overlapping fragments of the multi-agent coordination problem. Four broad
surveys already exist. What is missing is a principled basis for comparing
them.

{\bf Objectives:}
We set out to make the comparison auditable rather than narrative. That means
one rubric applied uniformly to every protocol, one recorded reason per
score, and an agenda of open problems that follows from the resulting gaps
rather than from opinion.

{\bf Methods:}
We contribute a uniform \emph{rubric} of ten evaluation axes. The axes span
discovery, identity, authorization, transport, task model, async, multi-modal
handling, security threat surface, semantic interoperability, and governance.
We survey thirteen protocols across three architectural layers (tool-binding,
peer-messaging, commerce-settlement). We score the eight tool-binding and
peer-messaging protocols against the full rubric. The five
commerce-settlement protocols are compared on a separate structural basis,
because the rubric's task, transport, and threat axes do not apply to a
settlement envelope. Scores come from a single rater reading each
specification under stated calibration rules. We report an auditable change
history over a machine-readable score file, and no inter-rater agreement
statistic.

{\bf Results:}
No protocol scores uniformly high, and the rubric shows why: the
contemporary protocols specialize into mutually exclusive strengths. None
combines a deployed task model with a decentralized identity model. The
authorization axis is the sharpest result. Its top level is empty across all
eight scored protocols, so no surveyed peer-messaging protocol offers a
standardized delegation credential that a downstream agent can verify. The
gap is universal rather than near-universal, and it holds even for the
protocol with the widest threat coverage in the set. We also map industry
adoption across enterprise SaaS, agentic commerce and payments, healthcare,
and finance, and we survey the emerging security and threat-modeling
literature.

{\bf Conclusions:}
The gaps identify the work of the next standardization cycle, and we propose
three co-designed protocol extensions against them. \emph{OWA} is a URI
scoping scheme. It adds an explicit \emph{workspace} tier between the
organizational trust root and the agent identifier. \emph{ACSP} is a sibling
envelope that records the sender's declared context-selection step: which
candidate items were considered, which were kept, which were dropped, under
which token budget and optimisation objective. Its signed wire-level audit
log carries observable invariants the receiver can verify. \emph{AGRP} is a
normative A2A envelope. It binds every cross-boundary message to a
classification-label set and a commitment-bearing redaction log on the spans
the egress filter claims to have touched. It also requires either a signed
Tier~1 enforcement claim or independently verified Tier~2 evidence. The three
close a single loop, naming to selection to attestation, that no contemporary
protocol provides. All three are positioned as specification-sketch
contributions for the next AAIF and W3C revision cycles rather than as
validated standards.
\end{abstract}

% JAIR Note from the AuthorKit: do not include ACM CCS Concepts or Keywords.
% The \keywords list that stood here under acmart was removed for that reason.

% Placeholder submission date. JAIR fills the accepted date at production;
% \received[accepted]{} is added then, not now. The arXiv posting omits it:
% a received date is the journal's record of a submission, and printing it on
% a preprint asserts a review process the preprint is not part of. arXiv
% stamps its own submission date. \ifarxiv is set in main.tex.
\ifarxiv\else
  \received{14 September 2026}
\fi

\maketitle

% Cutoff date declared on page 1, per the differentiation strategy (R1, R3).
\noindent\textbf{Note on scope.} Specifications and literature surveyed
up to 2026-Q1; later developments may be referenced in footnotes but
are not exhaustively covered.
